\chapter{Principi fondanti}
\label{chap:principi-fondanti}

\section{Zero-trust}
\label{sec:sota-zero-trust}

Il termine "zero trust" venne coniato nel 2010 dall'analista Forrester John Kindervag. L'idea di fondo 
era quella di non concedere fiducia in base alla posizione di rete, ma verificare ogni accesso
indipendentemente da essa.

Il NIST Special Publication 800-207 \cite{nist-sp-800-207} formalizza oggi questo principio come l'assunzione che nessuna rete,
interna o esterna, debba essere considerata implicitamente affidabile, e che ogni richiesta di accesso
vada valutata indipendentemente dalla sua provenienza.

RAM-USB estende il principio zero-trust evitando di credere persino all'identità di un componente 
autenticato via mTLS. 


\section{Defense-in-depth}
\label{sec:sota-defence-in-depth}

Il NIST Special Publication 800-53 Revision 5 definisce la defense-in-depth come una strategia che integra persone, tecnologia e capacità operative per 
stabilire barriere ridondanti su più livelli di un'organizzazione, in modo che il guasto o l'aggiramento di 
una singola barriera non comprometta il sistema nel suo complesso \cite{nist-sp-800-53r5}.